\section{1840(1256) Tarihli Ceza Kanunnamesi}
\indent\indent Osmanlı Devleti'nde bir ceza kanunnamesinin teşekkülü için 19.yüzyılı beklemek gerekecekti. 19. yüzyıl, devlet teşkilatlanması açısından 
modernizasyon ile anılmaktadır. Bürokrasi, devlet mekanizmasındaki sorunların farkına 16.yüzyıl gibi erken bir dönemde varmıştı. Bu sorunların çözümü için 
önerilen tavsiylerin en meşhuları Koçi Bey'in IV.Murad'a sunduğu risale ile Defterdar Sarı Mehmed Paşa'nın III.Ahmed'e sunduğu \textit{Nesâyihü'l-Vüzerâ ve'l-ümerâ} lahiyasıdır. 
III.Selim'in Nizâm-ı Cedîd teşebbüsüne kadar çözüm hep kânun-ı kadîm üzerinden aranmaya çalışılmıştır. Tanzimat Fermanı ile birlikte, ilmiye sınıfının çözüm olarak 
sunduğu kânun-ı kadîm anlayışı terkedilmiş ve kalemiye sınıfının ortaya koyduğu batı tarzı müesseseleşme yoluna girilmiştir. Bunun sonucunda Osmanlı bürokrasisinde bir ilmiye-kalemiye 
çatışması ortaya çıkmıştır. 
\subsection{Tanzimat Fermanı}
\indent\indent Gülhâne Hâtt-ı Şerîfi, Gülhâne Hâtt-ı Hümâyunu veya Tanzîmât-ı Hayriye olarak da bilinen Tanzimat Fermanı, 3 Kasım 1839 tarihinde Mustafa Reşid Paşa tarafından Gülhane Parkı'nda okunarak ilan edilmiştir.
Tanzimat Fermanı(1839), Osmanlı Devleti'nin modernleşme sürecinde önemli bir dönüm noktasıdır. Fermana giden süreci anlamak için öncesindeki 50 yıllık süreçte meydana gelen gelişmeleri anımsamak yardımcı olacaktır.
\subsubsection{Tanzimat Fermanı'na Giden Süreç}
III.Selim, 1787-92 Osmanlı-Rus Savaşı ve 1788-91 Osmanlı-Avusturya Savaşı 
sonrasında devlet müesseselerinde batı tarzı usulleri bensimeyen reformlar yapılmasında karar kıldı. Bu niyetle Nizâm-ı Cedîd hareketini başlattı. Adından da 
anlışalıcağı üzere, bu hareketteki amaç yeni bir düzen tesis etmekti. Uzun soluklu olmayan bu girişim, III.Selim'in 1807 yılında tahttan indirilmesi ile sonuçlandı.

III.Selim'in tahttan indirilmesi, Osmanlı Devleti'nde yaklaşık 30 sene sürecek bir otorite buhranına sebep oldu dersek abartmış olmayız. 
Ortaya çıkan krizin çözümüne yönelik ilk girişim Sened-i İttifak(1808) ile olmulştur. Sened-i İttifak, devlet otoritesininde sadece merkezin değil, ayanların da söz sahibi olacağı yönetime dair bir vesikadır. 
Vesikada, devlet bürokratları ve taşra hanedanları arasındaki mücadelenin devleti yıkıma sürüklediği belirtiliyor ve senet ile problemin çözümü için birlik şartlarını tesis ediliyor. Devletin dayandığı temel padişah otoritesidir 
ve bu otorite vüzera, ulema, rical, hanedanlar ve ocaklar birlikte çalıştığı sürece tesis edilebilir. Padişahın mutlak vekili olarak sadrazam kabul edilse de \textit{sadaret makamı da kanun dışı 
irtikâp ve irtişada bulunursa veya devlete zararlı kötü işlere kalkışırsa}, bu durum ittifak içinde önlenmelidir. Her ne kadar II.Mahmud bu misâkın dışında tutulsa da, padişahın mutlak otoritesini sınırlayıcı özelliği bakımından kayda değerdir.\footfullcite[92-95]{tanzimat}

Senede karşı ilk muhalefet, senedin mimarı olan Alemdar Mustafa Paşa'nın Sekban-ı Cedid Ocağı adlı askeri ıslahata giriştiğinde ortaya çıktı. Bu ıslahat karşısında yeniçeriler ayaklandı. Merkezi otorite ya yeniçerleri ya da
ayanlar ve ıslahatı destekleyecekti. Sened, merkezi otoriteyi zayıflattığı için padişah yeniçeriler tarafında saf aldı. Kasım 1808'de Alemdar Mustafa Paşa öldürülerek saf dışı bırakıldı.\footcite[139]{tanzimat}

Bu noktada, "Merkezi otorite iki ay içinde senedin mimarını öldürecekse, neden en başta kabul etti?" sorusu akıllara gelebilir. Sened, II.Mahmud'un cülusundan yaklaşık iki ay sonra, 1806-1812 Osmanlı-Rus Savaşı'nın devam ettiği esnada imzalanmıştır. Devlet, asker toplamak için 
ayanlara muhtaçtı. Fırsat buldukça ayanların gücü kırılmaya çalışıldıysa da, daha net teşebbüsler için savaşın bitişi olan 1812 yılını beklemek gerekcekti. Özellikle 1820 yılında, ayanların en güçlüsü Yanya Âyanı Tepedelenli Ali Paşa'nın üzerine yürünmesi ve 17 aylık mücadele 
sonucunda ele geçirilmesi ile en ciddi adım atılmış oldu. Merkezi otorite önündeki en ciddi engellerden biri kaldırılmış olsa da, bölgedeki Rumları isyandan yıldıran figürün ortadan kaldırılması başka sorunları beraberinde getirdi. Yanya kuşatmasından yedi ay sonra 
Rumlar ayaklandı. İsyanı bir türlü bastırılamayınca II.Mahmud Mısır Valisi Mehmed Ali Paşa'dan yardım istedi(1824).\footcite[140]{tanzimat} İbrahim Paşa'nın kısa sürede Mora'yı geri alması yeniçerilerin perişanlığını ortaya koyuyordu. Nitekim, 1826 yılında yeniçeri ocağı ortadan 
kaldırıldı. 

Öte yandan 1825 yılının sonuda vefat eden Rus Çarı I.Aleksandr'ın yerine I.Nikola'nın geçmesi isyanın seyrini değiştirdi. I.Nikola'nın daha agresif bir dış 
politika izlemeye niyetliydi. Bu politikanın bir sonucu olarak, 20 Ekim 1827'de Navarin'de Osmanlı donanması imha edildi. Sonrasında 1828-29 Osmanlı-Rus Savaşı'nda alınan mağlubiyet ile Yunanistan'ın bağımsızlığı tanındı. 

Yunan İsyanı, Osmanlı Devleti'nin iç isyanını bile bastıramayacak bir durumda olduğunu gözler önüne seriyordu. Mısır Valisi Kavalalı Mehmed Ali Paşa, belki bu durumdan aldığı 
cesaretle başka bir isyana sebep olacaktır. Kendi valisinin isyanını bastıramayan Osmanlı Devleti, 1833 Hünkâr İskelesi Antlaşması ile Rusya'nın desteğine başvuracaktır. Sular bir süreliğine durulusa da,
İbrahim Paşa komutasındaki Mısır ordusu, Osmanlı Devleti'ne Nizip'te(1839) küçük düşürücü bir bozgun yaşatacaktır.

Tanzimat Fermanı'nın yayınlandığı 1839 Kasımı'na gelirken, Osmanlı Devleti'nin merkezi otoritesi oldukça zayıflamış durumdaydı. Fermanın mimari Mustafa Reşid Paşa'ya göre çözüm, batı tarzı müesseselerin teşkil edilmesindeydi.

\subsubsection{Mustafa Reşid Paşa}
\indent\indent Bu noktada ismi Tanzimat Fermanı ile özdeşleşen Mustafa Reşid Paşa'dan bahsetmek gereklidir. Erken yaşlarda eniştesi Ispartalı Seyid Ali Paşa ile Osmanlı bürokrasisine adım atan Mustafa Reşid Paşa,
1834 yılında Paris sefiri ve daha sonra da Londra sefiri olarak görev yapmıştır.\footcite[76]{tarih_musahabeleri} Bu sefirlikleri sırasında Paris'te Thiers, Guizot ve Londra'da Palmerston gibi dönemin önemli devlet adamlarıyla tanışma fırsatı bulmuştur.\footcite[124]{tanzimat}

Mustafa Reşid Paşa'nın sefirliği esnasında Avrupa din merkezli devlet anlayışını terk edip, insan merkezli bir yapıya geçişin sancılarını yaşamaktaydı. Fransa'da \textit{Libertè, Ègalitè, Fraternitè(Hürriyet, Eşitlik, Kardeşlik)} sloganıyla 1789 yılında ortaya çıkan ihtilal, 
yönetme erkinin kaynağının Tanrı değil, halk olduğunu ilan ediyordu. Bu sayede halk teb`a olmaktan çıkıp vatandaş haline geliyordu. Bu sayede Orta Çağ'ın din temelli devlet anlayışı yıkılıp, modern devlet anlayışı ortaya çıkıyordu. 

Öte yandan II.Mahmud gibi inatçı, uzlaşmaya uzak bir padişah da Mustafa Reşid Paşa'nın fikirleri üzerinde etkili olmuştur. Özellikle genç yaşlarında maiyetinde bulunduğu Pertev Efendi'nin idam edilmesi,
Mustafa Reşid Paşa üzerinde derin etki bırakmıştı. Şahısların kanun önünde yargılanmadan, keyfi kararlarla cezalandırılmamaları gerekliliğini savunuyordu. Yine II.Mahmud'un Tepedelenli Ali Paşa'ya karşı yürüttüğü Yanya Kuşatması'nda, Mora'da çıkan Rum İsyanı'na rağmen 
devam etmekte ısrarcı olması siyasi olarak bir hata görülmekteydi. Mısır melesinde de Kavalılarla uzlaşmak yerine Hünkar İskelesi Anlaşması'nı Ruslarla imzalamak pahasına savaşa devam etmesi bir başka siyasi hata olarak yorumlanmıştı.\footcite[143]{tanzimat} 
İşte bu tecrübelerin etkisiyle Mustafa Reşid Paşa, devletin kurtuluşu için batı tarzı modernizasyonu gerekli görüyordu. Devlet yönetiminin şahıslara değil müesseselerei bağlı olması inancındaydı. Bundan dolayıdır ki, Tanzimat Fermanı'nın mimarı olması şaşırtıcı değildir.

\subsubsection{Tanzimat Fermanı'nın İçeriği}
\indent\indent Fermanın, devletin mevcut durumunun izahı ile başlamaktadır. Osmanlı Devleti'nin kuruluşundan beri şeri`ata bağlı kaldığını ve bu sayede yükseldiği fakat 150 yıldır çeşitli nedenlerden dolayı bu geleneğin terk edilderek devletin 
güçsüz ve fakir düştüğü belirtilmektedir. Gereken önlemler alındığı ve işlere girişildiği takdirde beş on yıl içinde bütün sorunların çözülebileceği, bunun için de:
\begin{itemize}
    \item Devletin ve memleketin idaresi için yeni kanunlar koyulması
    \item Can, mal, ırz ve namus emniyetinin devletin temel görevler olarak üstlenmesi
    \item Vergilerin gelir esasa alınarak daha adaletli toplanması 
    \item Gereken askerin toplanması ve görev süresi belierlenmesi
\end{itemize}
gereklilikleri vurgulanmıştır.\footcite[100]{tanzimat} 

Teb`a can, mal, ırz ve namusundan emin olursa bütün dikkat ve gayretini memleketin kalkınması ve refahı için sarf edeceğini belirtmektedir. Bu sebepten dolayı müslüman ve diğer milletlerin can, mal, ırz ve namuslarının emniyeti
istinasız olarak, Şeri`at hükümleri gereğince, devlet tarafından tam garanti edilmektedir.\footcite[101]{tanzimat} Ferman, bütün teb`aya şahsi emniyet ve ferdi haklar verildiği açıktır. Bu haklara karşı işlenen tecavüzlere matuf ceza 
ve usullerin berlinlenmesi gerekliliği doğmaktadır. Bu yüzdendir ki, Tanzimat Fermanı'ndan kısa bir süre sonra 1840 Ceza Kanunnamesi hazırlanmıştır.


\subsection{1840 Ceza Kanunnamesinin İçeriği}
\indent\indent 1840(1256) tarihli ceza kanunnamesi, mukaddime, on üç fasıl, kırk bir madde ve hatîmeden meydana gelmektedir. Mukaddimede açık bir şekilde Tanzimat Fermanı'na atıf bulunmkatadır. Mukaddime şöyle başlamaktır:
\begin{quote}
Beyaz üzerine şeref-sudür buyurulup geçen elli beş senesi Şa'ban-ı Şerifinin yirmi altıncı günü Gülhane'de kıraat olunan Hatt-ı Hümâyün-u adalet-makrün-i Hazret-i Şahâne mücebince kâffe-i tebea-i Devlet-i Aliyye bila istisna emniyet-i can ve mal ve mahfüziyet-i ırz ve namus hukuk-u mefrüzasına ez-ser nev-nail olmuş ve ber-muktezay-i hürriyet-i şer`iyye huzuru şer` ve kanunda ve mevadd-ı hukukiyede herkesin yeksan ve siyyan olması umür-i tabiiyyeden bulunmuş olmakla...\footfullcite[810-811]{islam_osmanli_hukuku} 
\end{quote}

Kanunnamedeki fasılları aşağıdaki şekilde özetleyebiliriz:\footcite[811-818]{islam_osmanli_hukuku}
\begin{itemize}
    \item Birinci Fasıl: Devlet aleyhine fitne ve katl suçları ve cezaları(kısas) üzerine dört madde
    \item İkinci Fasıl: Söz ve hareket ile karışıklık çıkarma suçu ve cezası(kürek) üzerine dört madde
    \item Üçüncü Fasıl: Irz ve namusa karşı işlenen suçlar(hakaret, darb vb.) ve cezaları(hapis) üzerine beş madde
    \item Dördüncü Fasıl: Mal ve emlâka karşı işlenen suçlar(hırsızlık, el koymak vb.) ve cezaları(sürgün) üzerine iki madde
    \item Beşinci Fasıl: Rüşvet suçu ve cezaları(tenzil-i rütbe, azil, kürek, sürgün) üzerine yedi madde
    \item Altıncı Fasıl: Devlet memurlarının işlediği hırsızlık suçu ve cezaları(tenzil-i rütbe, azil, kürek) üzerine üç madde
    \item Yedinci Fasıl: Devlet malının çalınması suçu ve cezaları(azil) üzerine dört madde
    \item Sekizinci Fasıl: Devlet görevlilerinin sorumluluklar ve anlaşmazlıklarında uygulanacak ceza(tedip) üzerine iki madde
    \item Dokuzuncu Fasıl: Verginin büyük küçük herkesin mâli kuvvetine göre toplanmasına, zaptiye ve raptiyeye muhalefet ve mukavemet suçu ve cezaları(hapis, kürek, kısas) üç madde
    \item Onuncu Fasıl: Silahlı yaralama ve öldürme suçu ve cezaları(kürek, kısas) üzerine iki madde
    \item On Birinci Fasıl: Yol kesme suçu ve cezaları(kürek, kısas) üzerine üç madde
    \item On İkinci Fasıl: Herkesin kanun önünde eşit olduğu ve kimsenin sebepsiz yere cezalandırılamayacağı üzerine iki madde
    \item On Üçüncü Fasıl: Taşra memurlarına muhalefet suçunun da işbu ceza kanunnamesinin kapsamına girmesi
\end{itemize}

Yukarıda da görüldüğü üzere, 1840 Ceza Kanunnamesi devlete ve toplum asayişine karşı işlenen suçları ve cezaları düzenlemiştir. Ceza türlerini genel olarak kısas, kürek, nefy(sürgün) ve hapis olarak dört ana kategoride toplamak mümkündür. 
Devlet memurlarınının işlediği suçlar tenzil-i rütbe, azil, tedip gibi daha hafif cezalar ile düzenlenmiştir. Hatîmede ise bir kere daha bu kanunnamenin \textit{Devlet-i Aliyyenin merâtib-i adide ve mütefavitesinde bulunanlar ve ehl-i islam ve milel-i saireden olanlar haklarında bila istisna vela garaz icra olunmak üzere karargir}
olduğu belirtilmiştir.\footcite[819]{islam_osmanli_hukuku}